{{دیوان
| عنوان = وطنیّت
| ماخذ = ویکی ماخذ
| ماخذ_لنک = https://ur.wikisource.org/wiki/%D9%88%D8%B7%D9%86%DB%8C%D9%91%D8%AA
| مقابلہ = اقبال سائبر لائبریری (اقبال اکادمی پاکستان)
| مقابلہ_لنک = https://www.iqbalcyberlibrary.net
}}
=== (یعنی وطن بحیثیّت ایک سیاسی تصوّر کے) ===

<poem>
اس دور میں مے اور ہے، جام اور ہے جم اور
ساقی نے بِنا کی روِشِ لُطف و ستم اور

مسلم نے بھی تعمیر کِیا اپنا حرم اور
تہذیب کے آزر نے ترشوائے صنم اور

ان تازہ خداؤں میں بڑا سب سے وطن ہے
جو پیرہن اس کا ہے، وہ مذہب کا کفن ہے

یہ بُت کہ تراشیدۂ تہذیبِ نوی ہے
غارت گرِ کاشانۂ دینِ نبَوی ہے

بازو ترا توحید کی قوّت سے قوی ہے
اسلام ترا دیس ہے، تُو مصطفوی ہے

نظّارۂ دیرینہ زمانے کو دِکھا دے
اے مصطفَوی خاک میں اس بُت کو ملا دے!

ہو قیدِ مقامی تو نتیجہ ہے تباہی
رہ بحر میں آزادِ وطن صُورتِ ماہی

ہے ترکِ وطن سُنّتِ محبوبؐ الٰہی
دے تُو بھی نبوّت کی صداقت پہ گواہی

گُفتارِ سیاست میں وطن اور ہی کچھ ہے
ارشادِ نبوّت میں وطن اور ہی کچھ ہے

اقوامِ جہاں میں ہے رقابت تو اسی سے
تسخیر ہے مقصودِ تجارت تو اسی سے

خالی ہے صداقت سے سیاست تو اسی سے
کمزور کا گھر ہوتا ہے غارت تو اسی سے

اقوام میں مخلوقِ خدا بٹتی ہے اس سے
قومیّتِ اسلام کی جڑ کٹتی ہے اس سے
</poem>
